\section{Equipes de Fórmula SAE}
A Fórmula SAE é uma competição estudantil de âmbito mundial, organizada no Brasil
pela SAE Brasil e realizada anualmente \cite{sae_brasil_formula_sae}. As equipes são compostas por alunos
de engenharia de uma mesma universidade, que ao longo do ano projetam,
constroem e testam um protótipo de veículo monoposto de competição. Na edição
de 2025, a competição brasileira contou com 66 equipes inscritas, representando
64 instituições brasileiras de ensino superior de 11 estados e do Distrito
Federal, além de duas equipes do exterior \cite{sae_brasil_formula_sae_2025_inscritos}.

A competição é estruturada em provas estáticas e dinâmicas. As provas estáticas
avaliam o projeto de engenharia, o plano de negócios e o detalhamento de custos
e manufatura; as provas dinâmicas testam o desempenho do protótipo em situações
como aceleração, skidpad, autocross e endurance. A prova de endurance tem maior
peso na pontuação e exige que o veículo percorra aproximadamente 22 km sem
falhas \cite{fsaeonline_endurance_2026}. Dado o calendário anual, cada ciclo de
desenvolvimento é único: falhas nas provas representam a perda de um ano de
trabalho, o que impõe requisitos de robustez, confiabilidade e melhoria
contínua sobre todos os subsistemas do veículo.

Entre as equipes que participam da Fórmula SAE, a EESC-USP Formula SAE é o
caso estudado neste trabalho. Vinculada à Escola de Engenharia de São Carlos
(EESC) da Universidade de São Paulo (USP), a equipe reúne estudantes de vários
cursos para projetar, manufaturar e testar protótipos em competições nacionais
e internacionais \cite{eesc_usp_formula_sae_e22_2025}. Seu título mais recente foi
conquistado em 2025, quando venceu a categoria Combustão da 21ª Competição
Nacional de Fórmula SAE Brasil, garantindo vaga para a etapa mundial
\cite{eesc_usp_formula_sae_2025,sae_brasil_resultados_combustao_2025}.

Como equipe vencedora e de alto desempenho, a EESC-USP depende fortemente de
uma cadeia de suporte técnico que transforma dados experimentais em decisões
práticas. Esse vínculo entre a coleta de dados no carro e as intervenções no
projeto é o foco das próximas subseções.

\section{A Importância da Aquisição de Dados}
Em ambientes de competição automotiva, compreender o comportamento do veículo
é condição essencial para extrair seu máximo desempenho. Sensores distribuídos
ao longo do protótipo fornecem informações que permitem validar projetos de
engenharia, orientar ajustes de setup, monitorar a integridade dos sistemas
críticos em tempo real e comparar objetivamente o desempenho entre pilotos
\cite{schultz2017_formula_sae_telemetry,ashford2020_formula_sae_data_analysis}. Sem
instrumentação adequada, a equipe fica limitada ao feedback subjetivo dos
pilotos — valioso, mas muitas vezes insuficiente para decisões precisas.

A aquisição de dados converte observações qualitativas em grandezas
mensuráveis: pressão de freio, curso de suspensão, ângulo de esterço,
aceleração lateral, temperatura de óleo, entre outras. Esses sinais permitem
identificar trechos da pista onde o carro perde tempo, avaliar diferenças entre
pilotos e verificar se modificações realizadas entre sessões surtem o efeito
esperado. Assim, o sistema de aquisição funciona como o elo entre o trabalho
de bancada e as decisões tomadas em pista. Trabalhos acadêmicos mostram
diversas soluções de baixo custo para telemetria e aquisição em Fórmula SAE,
com combinações variadas de armazenamento local, comunicação sem fio e
visualização em tempo real \cite{stoffman2004_low_cost_racing_daq,visconti2019_formula_sae_telemetry,zanetti2013_formula_sae_wireless,lauro2021_formula_sae_telemetria}.

No caso da EESC-USP, o sistema de aquisição de dados foi consolidado ao longo de
vários anos de desenvolvimento e testes, em continuidade a trabalhos anteriores
de eletrônica embarcada para Fórmula SAE \cite{blatt2020_aquisicao_formula_sae}.

A arquitetura adotada baseia-se em dois módulos de aquisição com
microcontroladores STM32, interligados por um barramento CAN a 1 Mbit/s, que
também conecta outros módulos eletrônicos do veículo (controle, transmissor
e distribuição de energia), formando uma rede de nós de comunicação
coesa \cite{silva2009_rede_can_formula_sae}. Cada módulo lê sinais provenientes
de sensores analógicos e digitais instalados no carro — acelerômetro, giroscópio,
sensores de pressão de freio, potenciômetros de suspensão e de esterço, entre
outros — e os transmite periodicamente pela rede CAN.

Os dados são registrados localmente em cartão SD embarcado. Ao final de cada
sessão, os arquivos são extraídos e processados em um software dedicado, que
gera arquivos estruturados para análise em ferramentas especializadas. Esse
fluxo (registro embarcado + pós-processamento) permite à equipe acompanhar a
evolução do veículo e dos pilotos ao longo do ano, fornecendo uma base objetiva
para decisões de engenharia e planejamento de melhorias.

